
%% bare_jrnl.tex
%% V1.4b
%% 2015/08/26
%% by Michael Shell
%% see http://www.michaelshell.org/
%% for current contact information.
%%
%% This is a skeleton file demonstrating the use of IEEEtran.cls
%% (requires IEEEtran.cls version 1.8b or later) with an IEEE
%% journal paper.
%%
%% Support sites:
%% http://www.michaelshell.org/tex/ieeetran/
%% http://www.ctan.org/pkg/ieeetran
%% and
%% http://www.ieee.org/

%%*************************************************************************
%% Legal Notice:
%% This code is offered as-is without any warranty either expressed or
%% implied; without even the implied warranty of MERCHANTABILITY or
%% FITNESS FOR A PARTICULAR PURPOSE! 
%% User assumes all risk.
%% In no event shall the IEEE or any contributor to this code be liable for
%% any damages or losses, including, but not limited to, incidental,
%% consequential, or any other damages, resulting from the use or misuse
%% of any information contained here.
%%
%% All comments are the opinions of their respective authors and are not
%% necessarily endorsed by the IEEE.
%%
%% This work is distributed under the LaTeX Project Public License (LPPL)
%% ( http://www.latex-project.org/ ) version 1.3, and may be freely used,
%% distributed and modified. A copy of the LPPL, version 1.3, is included
%% in the base LaTeX documentation of all distributions of LaTeX released
%% 2003/12/01 or later.
%% Retain all contribution notices and credits.
%% ** Modified files should be clearly indicated as such, including  **
%% ** renaming them and changing author support contact information. **
%%*************************************************************************


% *** Authors should verify (and, if needed, correct) their LaTeX system  ***
% *** with the testflow diagnostic prior to trusting their LaTeX platform ***
% *** with production work. The IEEE's font choices and paper sizes can   ***
% *** trigger bugs that do not appear when using other class files.       ***                          ***
% The testflow support page is at:
% http://www.michaelshell.org/tex/testflow/



\documentclass[journal]{IEEEtran}
%
% If IEEEtran.cls has not been installed into the LaTeX system files,
% manually specify the path to it like:
% \documentclass[journal]{../sty/IEEEtran}





% Some very useful LaTeX packages include:
% (uncomment the ones you want to load)


% *** MISC UTILITY PACKAGES ***
%
%\usepackage{ifpdf}
% Heiko Oberdiek's ifpdf.sty is very useful if you need conditional
% compilation based on whether the output is pdf or dvi.
% usage:
% \ifpdf
%   % pdf code
% \else
%   % dvi code
% \fi
% The latest version of ifpdf.sty can be obtained from:
% http://www.ctan.org/pkg/ifpdf
% Also, note that IEEEtran.cls V1.7 and later provides a builtin
% \ifCLASSINFOpdf conditional that works the same way.
% When switching from latex to pdflatex and vice-versa, the compiler may
% have to be run twice to clear warning/error messages.



\usepackage{graphicx}

\usepackage{textcomp}
% *** CITATION PACKAGES ***
%
%\usepackage{cite}
% cite.sty was written by Donald Arseneau
% V1.6 and later of IEEEtran pre-defines the format of the cite.sty package
% \cite{} output to follow that of the IEEE. Loading the cite package will
% result in citation numbers being automatically sorted and properly
% "compressed/ranged". e.g., [1], [9], [2], [7], [5], [6] without using
% cite.sty will become [1], [2], [5]--[7], [9] using cite.sty. cite.sty's
% \cite will automatically add leading space, if needed. Use cite.sty's
% noadjust option (cite.sty V3.8 and later) if you want to turn this off
% such as if a citation ever needs to be enclosed in parenthesis.
% cite.sty is already installed on most LaTeX systems. Be sure and use
% version 5.0 (2009-03-20) and later if using hyperref.sty.
% The latest version can be obtained at:
% http://www.ctan.org/pkg/cite
% The documentation is contained in the cite.sty file itself.






% *** GRAPHICS RELATED PACKAGES ***
%
\ifCLASSINFOpdf
  % \usepackage[pdftex]{graphicx}
  % declare the path(s) where your graphic files are
  % \graphicspath{{../pdf/}{../jpeg/}}
  % and their extensions so you won't have to specify these with
  % every instance of \includegraphics
  % \DeclareGraphicsExtensions{.pdf,.jpeg,.png}
\else
  % or other class option (dvipsone, dvipdf, if not using dvips). graphicx
  % will default to the driver specified in the system graphics.cfg if no
  % driver is specified.
  % \usepackage[dvips]{graphicx}
  % declare the path(s) where your graphic files are
  % \graphicspath{{../eps/}}
  % and their extensions so you won't have to specify these with
  % every instance of \includegraphics
  % \DeclareGraphicsExtensions{.eps}
\fi
% graphicx was written by David Carlisle and Sebastian Rahtz. It is
% required if you want graphics, photos, etc. graphicx.sty is already
% installed on most LaTeX systems. The latest version and documentation
% can be obtained at: 
% http://www.ctan.org/pkg/graphicx
% Another good source of documentation is "Using Imported Graphics in
% LaTeX2e" by Keith Reckdahl which can be found at:
% http://www.ctan.org/pkg/epslatex
%
% latex, and pdflatex in dvi mode, support graphics in encapsulated
% postscript (.eps) format. pdflatex in pdf mode supports graphics
% in .pdf, .jpeg, .png and .mps (metapost) formats. Users should ensure
% that all non-photo figures use a vector format (.eps, .pdf, .mps) and
% not a bitmapped formats (.jpeg, .png). The IEEE frowns on bitmapped formats
% which can result in "jaggedy"/blurry rendering of lines and letters as
% well as large increases in file sizes.
%
% You can find documentation about the pdfTeX application at:
% http://www.tug.org/applications/pdftex





% *** MATH PACKAGES ***
%
\usepackage{amsmath}
\DeclareMathOperator*{\argmin}{argmin}
\usepackage{amssymb}
\usepackage{color,soul}
\usepackage{multirow}
%\usepackage{ftnxtra}
%\usepackage{fnpos} %
%\usepackage{stfloats}
% A popular package from the American Mathematical Society that provides
% many useful and powerful commands for dealing with mathematics.
%
% Note that the amsmath package sets \interdisplaylinepenalty to 10000
% thus preventing page breaks from occurring within multiline equations. Use:
%\interdisplaylinepenalty=2500
% after loading amsmath to restore such page breaks as IEEEtran.cls normally
% does. amsmath.sty is already installed on most LaTeX systems. The latest
% version and documentation can be obtained at:
% http://www.ctan.org/pkg/amsmath





% *** SPECIALIZED LIST PACKAGES ***
%
%\usepackage{algorithmic}
% algorithmic.sty was written by Peter Williams and Rogerio Brito.
% This package provides an algorithmic environment fo describing algorithms.
% You can use the algorithmic environment in-text or within a figure
% environment to provide for a floating algorithm. Do NOT use the algorithm
% floating environment provided by algorithm.sty (by the same authors) or
% algorithm2e.sty (by Christophe Fiorio) as the IEEE does not use dedicated
% algorithm float types and packages that provide these will not provide
% correct IEEE style captions. The latest version and documentation of
% algorithmic.sty can be obtained at:
% http://www.ctan.org/pkg/algorithms
% Also of interest may be the (relatively newer and more customizable)
% algorithmicx.sty package by Szasz Janos:
% http://www.ctan.org/pkg/algorithmicx




% *** ALIGNMENT PACKAGES ***
%
%\usepackage{array}
% Frank Mittelbach's and David Carlisle's array.sty patches and improves
% the standard LaTeX2e array and tabular environments to provide better
% appearance and additional user controls. As the default LaTeX2e table
% generation code is lacking to the point of almost being broken with
% respect to the quality of the end results, all users are strongly
% advised to use an enhanced (at the very least that provided by array.sty)
% set of table tools. array.sty is already installed on most systems. The
% latest version and documentation can be obtained at:
% http://www.ctan.org/pkg/array


% IEEEtran contains the IEEEeqnarray family of commands that can be used to
% generate multiline equations as well as matrices, tables, etc., of high
% quality.




% *** SUBFIGURE PACKAGES ***
%\ifCLASSOPTIONcompsoc
%  \usepackage[caption=false,font=normalsize,labelfont=sf,textfont=sf]{subfig}
%\else
%  \usepackage[caption=false,font=footnotesize]{subfig}
%\fi
% subfig.sty, written by Steven Douglas Cochran, is the modern replacement
% for subfigure.sty, the latter of which is no longer maintained and is
% incompatible with some LaTeX packages including fixltx2e. However,
% subfig.sty requires and automatically loads Axel Sommerfeldt's caption.sty
% which will override IEEEtran.cls' handling of captions and this will result
% in non-IEEE style figure/table captions. To prevent this problem, be sure
% and invoke subfig.sty's "caption=false" package option (available since
% subfig.sty version 1.3, 2005/06/28) as this is will preserve IEEEtran.cls
% handling of captions.
% Note that the Computer Society format requires a larger sans serif font
% than the serif footnote size font used in traditional IEEE formatting
% and thus the need to invoke different subfig.sty package options depending
% on whether compsoc mode has been enabled.
%
% The latest version and documentation of subfig.sty can be obtained at:
% http://www.ctan.org/pkg/subfig




% *** FLOAT PACKAGES ***
%
%\usepackage{fixltx2e}
% fixltx2e, the successor to the earlier fix2col.sty, was written by
% Frank Mittelbach and David Carlisle. This package corrects a few problems
% in the LaTeX2e kernel, the most notable of which is that in current
% LaTeX2e releases, the ordering of single and double column floats is not
% guaranteed to be preserved. Thus, an unpatched LaTeX2e can allow a
% single column figure to be placed prior to an earlier double column
% figure.
% Be aware that LaTeX2e kernels dated 2015 and later have fixltx2e.sty's
% corrections already built into the system in which case a warning will
% be issued if an attempt is made to load fixltx2e.sty as it is no longer
% needed.
% The latest version and documentation can be found at:
% http://www.ctan.org/pkg/fixltx2e


%\usepackage{stfloats}
% stfloats.sty was written by Sigitas Tolusis. This package gives LaTeX2e
% the ability to do double column floats at the bottom of the page as well
% as the top. (e.g., "\begin{figure*}[!b]" is not normally possible in
% LaTeX2e). It also provides a command:
%\fnbelowfloat
% to enable the placement of footnotes below bottom floats (the standard
% LaTeX2e kernel puts them above bottom floats). This is an invasive package
% which rewrites many portions of the LaTeX2e float routines. It may not work
% with other packages that modify the LaTeX2e float routines. The latest
% version and documentation can be obtained at:
% http://www.ctan.org/pkg/stfloats
% Do not use the stfloats baselinefloat ability as the IEEE does not allow
% \baselineskip to stretch. Authors submitting work to the IEEE should note
% that the IEEE rarely uses double column equations and that authors should try
% to avoid such use. Do not be tempted to use the cuted.sty or midfloat.sty
% packages (also by Sigitas Tolusis) as the IEEE does not format its papers in
% such ways.
% Do not attempt to use stfloats with fixltx2e as they are incompatible.
% Instead, use Morten Hogholm'a dblfloatfix which combines the features
% of both fixltx2e and stfloats:
%
% \usepackage{dblfloatfix}
% The latest version can be found at:
% http://www.ctan.org/pkg/dblfloatfix




%\ifCLASSOPTIONcaptionsoff
%  \usepackage[nomarkers]{endfloat}
% \let\MYoriglatexcaption\caption
% \renewcommand{\caption}[2][\relax]{\MYoriglatexcaption[#2]{#2}}
%\fi
% endfloat.sty was written by James Darrell McCauley, Jeff Goldberg and 
% Axel Sommerfeldt. This package may be useful when used in conjunction with 
% IEEEtran.cls'  captionsoff option. Some IEEE journals/societies require that
% submissions have lists of figures/tables at the end of the paper and that
% figures/tables without any captions are placed on a page by themselves at
% the end of the document. If needed, the draftcls IEEEtran class option or
% \CLASSINPUTbaselinestretch interface can be used to increase the line
% spacing as well. Be sure and use the nomarkers option of endfloat to
% prevent endfloat from "marking" where the figures would have been placed
% in the text. The two hack lines of code above are a slight modification of
% that suggested by in the endfloat docs (section 8.4.1) to ensure that
% the full captions always appear in the list of figures/tables - even if
% the user used the short optional argument of \caption[]{}.
% IEEE papers do not typically make use of \caption[]'s optional argument,
% so this should not be an issue. A similar trick can be used to disable
% captions of packages such as subfig.sty that lack options to turn off
% the subcaptions:
% For subfig.sty:
% \let\MYorigsubfloat\subfloat
% \renewcommand{\subfloat}[2][\relax]{\MYorigsubfloat[]{#2}}
% However, the above trick will not work if both optional arguments of
% the \subfloat command are used. Furthermore, there needs to be a
% description of each subfigure *somewhere* and endfloat does not add
% subfigure captions to its list of figures. Thus, the best approach is to
% avoid the use of subfigure captions (many IEEE journals avoid them anyway)
% and instead reference/explain all the subfigures within the main caption.
% The latest version of endfloat.sty and its documentation can obtained at:
% http://www.ctan.org/pkg/endfloat
%
% The IEEEtran \ifCLASSOPTIONcaptionsoff conditional can also be used
% later in the document, say, to conditionally put the References on a 
% page by themselves.




% *** PDF, URL AND HYPERLINK PACKAGES ***
%
%\usepackage{url}
% url.sty was written by Donald Arseneau. It provides better support for
% handling and breaking URLs. url.sty is already installed on most LaTeX
% systems. The latest version and documentation can be obtained at:
% http://www.ctan.org/pkg/url
% Basically, \url{my_url_here}.







% correct bad hyphenation here
\hyphenation{}


\begin{document}
%
\title{Geometric Solution of Image Degradation by Diffraction in Lensless Sensing and Microscopy}
%\title{High Numerical Aperture Illumination for Non-Computational Lensless Sensing and Microscopy}
%
%

%

\author{Sanjeev~Kumar,~\IEEEmembership{}
        Manjunatha~Mahadevappa,~\IEEEmembership{Senior~Member,~IEEE,}
        and~Pranab~Kumar~Dutta,~\IEEEmembership{Member,~IEEE}% <-this % stops a space
\thanks{S. Kumar and M. Mahadevappa are with the School of Medical Science and Technology, Indian Institute of Technology Kharagpur, 721302, India, e-mail: sanjeevsmst@iitkgp.ac.in.}% <-this % stops a space
\thanks{P.K. Dutta is with the Department of Electrical Engineering, Indian Institute of Technology Kharagpur, 721302, India.}% <-this % stops a space



\thanks{Manuscript received date; revised date.}}





\markboth{Journal,~Vol, No.}%
{Kumar \MakeLowercase{\textit{et al.}}: Bare Demo of IEEEtran.cls for IEEE Journals}




% make the title area
\maketitle

% As a general rule, do not put math, special symbols or citations
% in the abstract or keywords.
\begin{abstract}
This paper proposes a non-computational method of counteracting the effect of image degradation introduced by the diffraction phenomenon in lensless microscopy. All the optical images (whether focused by lenses or not) are diffraction patterns, which preserve the visual information upto a certain extent determined by the size of the point spread functions, like airy disks in some cases. A highly diverging beam can be exploited to reduce the spatial extent of these point spread functions relatively in the transformed projective space, which can help us in the spatial unmixing of the visual information. The principle has been experimentally validated by the lensless imaging of red blood cells of diameter $\sim$6-9 $\mu$m and a photolithography mask with features in micrometer scale. The important advantages of the proposed approach of non-computational shadow microscopy are the improved depth of field and a drastic increase in the sensor to sample working distance. The imaging method can also be used as a projection technique in the multi-angle optical computed tomography (CT).
\end{abstract}

% Note that keywords are not normally used for peerreview papers.
\begin{IEEEkeywords}
Image degradation, Diffraction, Illumination, Sensing, Microscopy, Optical Fiber.
\end{IEEEkeywords}





\IEEEpeerreviewmaketitle

\section{Introduction}
\begin{figure}
\centering
{\includegraphics[width=6cm]{"figure1"}}
\caption{(a) A parallel beam illuminating two points and the corresponding image. (b) A diverging beam illuminating same two points and the corresponding image.}
\label{Fig:1}
\end{figure}
A lensless shadow microscopy system contains a light source to illuminate a weakly absorbing object and an image sensor to record this object\textquotesingle s shadow. Shadow microscopy with visible light has been discussed in the literature, for instance by \cite{yang2018resolution,fang2017chip,ozcan2016lensless} and \cite{wu2012optical}. The restriction to deposit sample directly on the sensor\textquotesingle s surface, for preserving the high resolution limited its use to on-chip cell or tissue cultures only. For instance, in the paper of Yang et al, 2018, the sample to sensor distance is 5 \ $\mu$m only, which is the thickness of the protective coating on the sensor surface \cite{yang2018resolution}. The fundamental problem that underlies lensless imaging is the scattering of the incident light by the sample and because of this phenomenon, every point in the object produces a cone of light which is projected as a point spread function (PSF) on the sensor. In incoherent illumination, this PSF can be approximated as a space-invariant Gaussian pattern whose variance depends on the sample to sensor distance \cite{muhammad2012sampling}. Because of the very narrow frequency bandwidth support of this Gaussian pattern, the reconstruction problem is severely ill-posed \cite{bertero2020introduction} and high resolution reconstruction of the object is very difficult (and has never been reported satisfactorily as per authors\textquotesingle \ literature survey). One solution is to use a coded mask in the beam path, so that the PSF takes the form of the magnified image of the mask itself (see \cite{wu2020single,cieslak2016coded,antipa2018diffusercam}). Now by the choice of a mask of larger frequency bandwidth support, a high resolution reconstruction of the object can be obtained using an appropriate algorithm. An alternate solution is to use a coherent illumination, where the PSF naturally has a very large frequency bandwidth support (unlike a Gaussian PSF). This approach is very common in lensless microscopy methods and is very well-known by the name of “lensless or digital in-line holographic microscopy” sometimes shortened as LIHM or DIHM (see \cite{ozcan2016lensless,kumar2020compressive,Kumar_2021,Kumar_2020} for details). Another similar lensless computational imaging technique used with the X-ray, ultraviolet light and electron beam is "coherent diffractive imaging" \cite{miao2011coherent,zurch2014real,clark2012high,yang2013coherent}.


\begin{figure}
\centering
{\includegraphics[width=8cm]{"figure7"}}
\caption{(a) A 512 $\times$ 512 sub-region from a 4096 $\times$ 4096 simulated test object. (b) Sub-region of the simulated diffraction image for the parallel beam illumination, sampled at pixel pitch of 0.125 $\mu$m, $z_2$ is the sample to sensor distance. (c) Simulated diffraction image for the Gaussian beam illumination corresponding to optical magnification value 11, downsampled to pixel pitch of 16 $\mu$m after numerical wave propagation. (d) Line plot of a single row (central row) from a and b. (e) Line plot of a single row (central row) from a and c, diffraction image is scaled appropriately before visualization to match the dimension of object in a.}
\label{Fig:1}
\end{figure}
\begin{figure*}
\centering
{\includegraphics[width=14cm]{"figure2"}}
\caption{Imaging setup for the proposed method.}
\label{Fig:2}
\end{figure*}

\section{Principle}
Unlike coded aperture imaging and LIHM, this paper does not consider solving the inverse problem for reconstruction of the object from the measurements. Instead this paper focuses on a geometric solution for preserving the information about the microscopic sample. A diffraction pattern is an image of the sample (and vice versa \cite{mertz2019introduction}) which contains the visual information upto a certain resolution, determined by the spatial extent of the point spread function (PSF). Suppose there are two points in the object which are illuminated with a parallel beam of light. For a certain sample to sensor distance $z_2$, suppose the PSF has a radius $d$. We assume the spatial invariance of this PSF in our field of view for the sake of simplicity in this discussion. Now the signals from the individual points will start overlapping if the distance between these two points becomes less than $2d$ (see figure 1a). We can safely comment here that the half-pitch resolution limit is the radius of the PSF itself. (We are not considering the Rayleigh limit because the PSF here is an arbitrary pattern which will depend on the coherence of the light and the strength of unscattered part of the light. PSF shown in figure 1 is only indicative). Alternatively, if a high numerical aperture (NA) illumination i.e. a diverging beam is used to illuminate the same two points, a magnification $M$ will be introduced in the image. In this case, the half pitch resolution limit is not the radius of the PSF ($d+ \delta d$) but the ratio of ($d+ \delta d$) and $M$. In other words, the centers of the patterns from the two individual points are obtained at a larger separation, due to the different geometric projection angles for the different points (see figure 1b). The value of magnification is the ratio of distance between the light source and sensor plane ($z$) and the distance between the light source and sample/object plane ($z_1$) i.e. $M=\frac{z}{z_1}$ \cite{forsyth2002computer,hartley2003multiple}. This geometric magnification can be effectively exploited for non-computational shadow microscopy only if the magnification is large enough to counteract the loss of resolution due to the degradation introduced by the PSF.
\begin{figure}
\centering
{\includegraphics[width=8cm]{"figure3"}}
\caption{(a-b) Projections (diffraction patterns) of red blood cells (RBCs) of diameter $\sim$6-9 $\mu$m at magnification $\sim$1, wavelength 670 nm. (c-d) Corresponding reconstructions using the angular spectrum method. All the four images are of digital resolution 256$\times$256 and pixel-pitch 1.12 $\mu$m.}
\label{Fig:3}
\end{figure}

\begin{figure*}
\centering
{\includegraphics[width=16cm]{"figure4"}}
\caption{Projections of RBCs using the proposed method at different magnifications, pixel-pitch is $\sim$18 $\mu$m for all the images corresponding to the proposed method. (a-c) At magnification $\sim$100 under three different illumination wavelengths. (d)  At magnification $\sim$150, wavelength 830 nm. (e-f) At magnification $\sim$200, wavelength 830 nm. (g-h) At magnification $\sim$200 with pseudocolor, wavelength 830 nm. (i) At magnification $\sim$100 visualized using pseudocolor and mesh plot, blue denotes high intensity and red denotes low intensity, wavelength 830 nm. (j) Under bright-field microscope with 40x objective lens and broadband illumination. (k) Under atomic force microscope.}
\label{Fig:4}
\end{figure*}

\begin{figure}
\centering
{\includegraphics[width=6cm]{"figure5"}}
\caption{(a) RBC under a bright field microscope with 100x objective lens and broadband illumination. (b) RBCs under an atomic force microscope showing the height using color.}
\label{Fig:5}
\end{figure}
\begin{figure}
\centering
{\includegraphics[width=9cm]{"figure6"}}
\caption{(a) A photolithography mask under a bright field microscope with 20x objective lens and broadband illumination. (b) Same mask\textquotesingle s image captured using the proposed method, digital resolution 178$\times$160, pixel-pitch $\sim$18 $\mu$m and wavelength 830 nm.}
\label{Fig:6}
\end{figure}

\section{Simulation experiments}
In this section, we validate the above described principle using the Fourier optics based simulation experiments.  For non-paraxial regime, the Rayleigh-Somerfeld (RS) diffraction integral or Angular spectrum (AS) method can be used for the numerical wave propagation \cite{goodman2005introduction}. AS method is computationally efficient and involves a point-wise multiplication with an appropriate optical transfer function in the frequency domain. The Fourier and the inverse Fourier transforms are obtained using the fast Fourier transform (FFT) and the inverse fast Fourier transform (IFFT) algorithms. The axial distance dependent optical transfer function $H(\textbf{v})$ can be obtained using the following equation \cite{kumar2020compressive,ozcan2016lensless}:
\begin{align}
H(\textbf{v}) &=  exp\Big(jk_0z_2 \sqrt{1 - (\lambda v_x)^2-(\lambda v_y)^2}\Big);  \sqrt{v^2_x+v^2_y} < \frac{1}{\lambda} \\
&= 0;\ \textrm{otherwise}
\end{align}
where $k_0 = \frac{2\pi}{\lambda}$ is the wave-number, $\lambda$ is the wavelength, $z_2$ is the distance, $(\textbf{v}=v_x,v_y)$ is the frequency coordinate vector.
Figure 2a shows a simulated test object with two bars of width 5 $\mu$m each. This image is a sub-region of a larger image of digital resolution 4096 $\times$ 4096 with a sampling distance of 125 nm. Figure 2b shows a sub-region of the simulated diffraction pattern for the case of parallel beam illumination, for a sample to sensor plane distance ($z_2$) of 100 $\mu$m. Figure 2b has the same digital resolution and spatial sampling as figure 2a. 


In the second experiment, the same test object is multiplied with a simulated Gaussian beam (diverging beam illumination) and then numerically propagated to the same distance $z_2$ = 100 $\mu$m. The same digital resolution of 4096 $\times$ 4096 and spatial sampling of 125 nm has been maintained. After numerical propagation, the obtained diffraction pattern is downsampled to a sampling distance of the 16 $\mu$m, to mimic the low resolution imaging sensor used in the bench-top experiments performed in the following sections of this paper. Figure 2c shows the full field of view of the obtained diffraction image. The optical magnification value corresponding to the figure 2c is 11. As described in the previous section, the visual information about the bars which is lost in the case of the parallel beam illumination is well preserved when a highly diverging beam illumination is used, for the same sample to sensor plane distance.
This numerical experiment has been restricted to a small sample to sensor plane distance ($z_2$) because of the following computational limitations: since a high numerical aperture beam has been simulated, the object must be zero padded upto the appropriate lateral lengths before numerical propagation step. As the value of $z_2$ scales up, the memory requirement increases, reaching finally upto the computer capacity. Computational complexity will also scale up for the same reason. Breaking down the problem into the propagation-downsampling-propagation loop will introduce numerical (aliasing) errors related to the downsampling. Tomasz Kozacki et al, 2012 and Tomasz Kozacki, 2008 described these sampling related numerical errors in the angular spectrum method \cite{kozacki2012computation,kozacki2008numerical}.



\section{Materials and methods}
A solid core photonic crystal fiber (PCF) of numerical aperture (NA) $0.38 \pm 0.05$ and an effective mode field diameter of $1.8 \pm 0.3 \ \mu m$ (both the values at 780 nm respectively) has been used as an illumination source in this paper. This NA value corresponds to a half cone angle of 22.33 degrees. For a step index single mode optical fiber, the NA value is around 0.1, which corresponds to a half cone angle of 5.74 degrees only. The fiber is connected (using FC/PC connector) to a pigtailed laser diode of wavelength 830 nm and power 10 mW but the power output is manually controlled using a diode controller while imaging, to obtain the optimum contrast. Some images have also been captured with illumination wavelengths 670 nm and 1310 nm for the generalization. Sample fixed on a glass slide is mounted on a micrometer stage to control its three-dimensional position precisely. This provides a mechanical control of both the magnification and the lateral scanning of the sample. A pinhole aperture of arbitrary diameter of $<$ 1 mm has been kept between the sample and the sensor to select the rays from the features of interest (to some extent). This helps us to filter out any rays from the high scattering angles from the features outside the field of view of interest.



In lensless in-line holographic microscopy (LIHM) at unit magnification, the pixel-pitch of the image sensor determines the limit of resolution (along with some other factors), unless some sub-pixel super-resolution technique is employed. So a high resolution image sensor is an essential requirement in LIHM. In the proposed method, the high resolution of the image sensor (i.e. small pixel-pitch) is of little interest but a significantly large sensing area is the essential requirement. The reason can be understood with this instance: if a 25 $\mu$m feature is magnified by 200 times, it will be projected on a 5 mm sensor area. In this paper, a lead-oxysulfide vidicon image sensor of 9.5 mm (vertical) $\times$ 12.7 mm (horizontal) sensing area has been used. The horizontal and vertical resolutions of this image sensor is around 18 $\mu$m both (9.5 mm/525 vertical raster scans and  12.7 mm/700 horizontal TV lines). 

\section{Results and discussions}
To experimentally demonstrate the principle discussed in this paper, we first show two diffraction patterns captured at magnifications around 1 in the figures 4a and 4b. These figures show the projections of isolated red blood cells (RBCs) of diameter $\sim$6-9 $\mu$m. For these two figures, the spatial extent of the PSFs is much larger than the size of object under observation and hence no visual information can be obtained directly by looking at these projections. Reconstructions from these diffraction patterns (using angular spectrum method \cite{matsushima2009band} based on the principle of in-line holography) have been included in figure 4c and 4d, to assist the reader in recognizing the degradations introduced in these images by the diffraction phenomenon. Next we introduce a magnification of around 100 times using the optical setup and imaging geometry shown in figure 3 and the results have been shown in figures 5a to c. Unlike the figure 4, isolated RBCs can be observed easily in these images. This change in the visual appearance of the diffraction pattern and the spatial unmixing of the signals from individual RBCs in the latter images is the direct experimental validation of the principle presented in this paper. Figures 5e to h correspond to an optical magnification of around 200 times obtained using the same imaging geometry, now a single RBC image extends to a length of around $\sim$1.5-2 mm on the image sensor. For instance, for the figure 5e, the digital resolution is 96$\times$104 and pixel pitch is $\sim$18 $\mu$m (pixel pitch is same for all the images acquired with setup shown in figure 3). In these images, even the well-known concave shaped morphology of the RBCs \cite{merola2017tomographic} can be undoubtedly observed. The light source and the image sensor\textquotesingle s positions are fixed, only the sample\textquotesingle s position is changed to control the magnification. As the magnification is increased, the field of view decreases proportionally. Figure 6 shows the diameter and the thickness of an RBC measured in a bright field microscope and an atomic force microscope as the gold standard methods. 

In figure 7, images of a photolithography mask with features of dimensions $\sim$600 nm (first bright vertical line like feature from left), $\sim$3 $\mu$m (next dark vertical feature) and $\sim$1 $\mu$m (second bright vertical line like feature) have been shown for the further validation of the imaging principle and resolution. A full pitch resolution of around $\sim$2-3 micrometers can be anticipated for the proposed method, from the images of the red blood cells and this photolithography mask.


The images captured using this imaging geometry have one-point perspective, due to the depth dependent magnification. One advantage of this method is that the depth of field is not limited to few microns (or less than a micron) like in lens based microscopy operated at similar magnifications. Also the sensor to sample distance is not restricted to a few microns like the previous demonstrations of the lensless shadow microscopy \cite{yang2018resolution}, this eliminates the restriction to deposit sample on the sensor surface. The optical fiber tip was in a close proximity of the sample ($\sim$50-200 $\mu$m) which can be handled using an appropriate instrumentation. The stochastic vibration of the cleaved end of the optical fiber tip is also a factor contributing to the blurring of images. The proposed method is real-time and free of any numerical or aliasing error because of its non-computational nature. At last, the utility of this principle can be found in the design and development of cell counters, flow cytometry, imaging in microfluidics, crack detection and in the development of novel microscopy technologies. The imaging method can also be used as a projection technique in the multi-angle optical computed tomography (CT).

\section{Conclusions}
In conclusion, the image degradation introduced in the lensless shadow microscopy by the diffraction phenomenon can be suppressed by changing the illumination strategy. Geometric magnifications of large values obtained using a large numerical aperture light source, can help to us perform lensless imaging without involving any computational reconstruction step. Sub-pixel resolution has been clearly demonstrated, as the pixel-pitch of the image sensor used was $\sim$18 $\mu$m and the sample being captured was of size $\sim$6-9 $\mu$m or of further smaller size. Working distance increased drastically from few micrometers to few millimeters using the proposed principle.






















% Can use something like this to put references on a page
% by themselves when using endfloat and the captionsoff option.
\ifCLASSOPTIONcaptionsoff
  \newpage
\fi



% trigger a \newpage just before the given reference
% number - used to balance the columns on the last page
% adjust value as needed - may need to be readjusted if
% the document is modified later
%\IEEEtriggeratref{8}
% The "triggered" command can be changed if desired:
%\IEEEtriggercmd{\enlargethispage{-5in}}

% references section

% can use a bibliography generated by BibTeX as a .bbl file
% BibTeX documentation can be easily obtained at:
% http://mirror.ctan.org/biblio/bibtex/contrib/doc/
% The IEEEtran BibTeX style support page is at:
% http://www.michaelshell.org/tex/ieeetran/bibtex/
%\bibliographystyle{IEEEtran}
% argument is your BibTeX string definitions and bibliography database(s)
%\bibliography{IEEEabrv,../bib/paper}
%
% <OR> manually copy in the resultant .bbl file
% set second argument of \begin to the number of references
% (used to reserve space for the reference number labels box)
\bibliographystyle{IEEEtran}
\bibliography{Sanjeev_NCLM}

%\begin{thebibliography}{1}

%\bibitem{IEEEhowto:kopka}
%H.~Kopka and P.~W. Daly, \emph{A Guide to \LaTeX}, 3rd~ed.\hskip 1em plus
%  0.5em minus 0.4em\relax Harlow, England: Addison-Wesley, 1999.

%\end{thebibliography}

% biography section
% 
% If you have an EPS/PDF photo (graphicx package needed) extra braces are
% needed around the contents of the optional argument to biography to prevent
% the LaTeX parser from getting confused when it sees the complicated
% \includegraphics command within an optional argument. (You could create
% your own custom macro containing the \includegraphics command to make things
% simpler here.)
%\begin{IEEEbiography}[{\includegraphics[width=1in,height=1.25in,clip,keepaspectratio]{mshell}}]{Michael Shell}
% or if you just want to reserve a space for a photo:

%\begin{IEEEbiography}{Michael Shell}
%Biography text here.
%\end{IEEEbiography}

% if you will not have a photo at all:
%\begin{IEEEbiographynophoto}{Sanjeev Kumar}




%\end{IEEEbiographynophoto}

% insert where needed to balance the two columns on the last page with
% biographies
%\newpage

%\begin{IEEEbiographynophoto}{Jane Doe}
%Biography text here.
%\end{IEEEbiographynophoto}

% You can push biographies down or up by placing
% a \vfill before or after them. The appropriate
% use of \vfill depends on what kind of text is
% on the last page and whether or not the columns
% are being equalized.

%\vfill

% Can be used to pull up biographies so that the bottom of the last one
% is flush with the other column.
%\enlargethispage{-5in}



% that's all folks
\end{document}


